\section*{Capítulo \ref{ch:b1:taylor} --- Fórmulas de Taylor e Expansões
Assintóticas}

\begin{solution}{exo:b1:taylor:1}
\[
\eu^{2x} = 1 + 2x + 2x^2 + \frac{4x^3}{3} + o(x^3);
\qquad
\ln(1 - x) = -x - \frac{x^2}{2} - \frac{x^3}{3} - \frac{x^4}{4} +
o(x^4);
\]
\[
\sqrt{1+x} = 1 + \frac x2 - \frac{x^2}{8} + \frac{x^3}{16} + o(x^3);
\qquad
\frac{1}{1 + x^2} = 1 - x^2 + x^4 - x^6 + o(x^6),
\]
a última substituindo $-x^2$ na expansão geométrica.
\end{solution}

\begin{solution}{exo:b1:taylor:2}
$\eu^x - 1 - x = \frac{x^2}{2} + o(x^2)$: limite $\dfrac12$.

$\cos x = 1 - \frac{x^2}{2} + \frac{x^4}{24} + o(x^4)$ e
$\sqrt{1 - x^2} = 1 - \frac{x^2}{2} - \frac{x^4}{8} + o(x^4)$:
diferença $\frac{x^4}{24} + \frac{x^4}{8} = \frac{x^4}{6} + o(x^4)$:
limite $\dfrac16$.

$\ln(1+x) - \sin x = \bigl(x - \frac{x^2}{2}\bigr) - x + o(x^2) =
-\frac{x^2}{2} + o(x^2)$: limite $-\dfrac12$.
\end{solution}

\begin{solution}{exo:b1:taylor:3}
$\frac{1}{1+t^2} = 1 - t^2 + t^4 + o(t^5)$; integrando de $0$ a
$x$ (\cref{met:b1:taylor:compute}~(5)):
\[
\arctan x = x - \frac{x^3}{3} + \frac{x^5}{5} + o(x^5)\ \ (\text{até mesmo }o(x^6)\text{, por imparidade}).
\]
$(1 - t^2)^{-1/2} = 1 + \frac{t^2}{2} + \frac38 t^4 + o(t^5)$ (expansão
binomial com $\alpha = -\frac12$, $x = -t^2$:
$\binom{-1/2}{2} = \frac{(-\frac12)(-\frac32)}{2} = \frac38$);
integrando:
\[
\arcsin x = x + \frac{x^3}{6} + \frac{3x^5}{40} + o(x^5) .
\]
\end{solution}

\begin{solution}{exo:b1:taylor:4}
Taylor--Lagrange (\cref{thm:b1:taylor:lagrange}) para $\exp$ em $a =
0$, $x = 1$: a \omterm{def:b1:derivative:def}{derivada} de ordem $(n+1)$ é $\eu^t \leq \eu < 3$ em
$\intcc{0}{1}$, logo
\[
\Bigl|\eu - \sum_{k=0}^{n} \frac{1}{k!}\Bigr| \leq
\frac{3}{(n+1)!} .
\]
Para $6$ decimais exatas, queremos $\frac{3}{(n+1)!} < 5\times 10^{-7}$,
isto é, $(n+1)! > 6\times 10^{6}$: como $10! = 3\,628\,800$ e $11! =
39\,916\,800$, temos $n + 1 = 11$, ou seja, $n = 10$ basta.
\end{solution}

\begin{solution}{exo:b1:taylor:5}
$n\ln\bigl(1 + \frac1n\bigr) = n\Bigl(\frac1n - \frac{1}{2n^2} +
\frac{1}{3n^3} + o\bigl(\frac{1}{n^3}\bigr)\Bigr) = 1 - \frac{1}{2n}
+ \frac{1}{3n^2} + o\bigl(\frac{1}{n^2}\bigr)$. Exponenciando, com
$u = -\frac{1}{2n} + \frac{1}{3n^2}$ e $\eu^u = 1 + u + \frac{u^2}2
+ o(u^2)$:
\[
\Bigl(1 + \frac1n\Bigr)^n = \eu\cdot \eu^{u}
= \eu\Bigl(1 - \frac{1}{2n} + \frac{1}{3n^2} + \frac{1}{8n^2}
+ o\Bigl(\frac{1}{n^2}\Bigr)\Bigr)
= \eu\Bigl(1 - \frac{1}{2n} + \frac{11}{24n^2} +
o\Bigl(\frac{1}{n^2}\Bigr)\Bigr).
\]
Limite $\eu$; o erro é $\sim \dfrac{\eu}{2n}$: lento (um dígito por multiplicação de
$n$ por dez).
\end{solution}

\begin{solution}{exo:b1:taylor:6}
$f(x) = x^2 - x^4 = x^2(1 + o(1))$: primeiro termo $x^2$, $p = 2$ par,
coeficiente $> 0$: mínimo local em $0$ (não global: $f(2) = -12$).

$g(x) = x^3 + x^5$: primeiro termo $x^3$, $p$ ímpar: nenhum extremo; $g$
atravessa a sua tangente (horizontal): inflexão em $0$.

$h = \exp$ em $a = 1$: $h(x) = \eu + \eu(x-1) + \frac{\eu}{2}(x-1)^2
+ o((x-1)^2)$; a diferença com a tangente é
$\frac{\eu}{2}(x-1)^2 + o(\cdot) > 0$ localmente. Globalmente: $\eu^x -
\eu x \geq 0$ para todo $x$, por \omterm{def:b1:derivative:convex}{convexidade}
(\cref{thm:b1:derivative:convexchar}~(3)): o gráfico está acima de toda
tangente, com igualdade apenas no ponto de contato.
\end{solution}

\begin{solution}{exo:b1:taylor:7}
Para $x \to +\infty$:
\[
f(x) = x\Bigl(1 + \frac1x\Bigr)^{1/3}
= x\Bigl(1 + \frac{1}{3x} - \frac{1}{9x^2} +
o\Bigl(\frac{1}{x^2}\Bigr)\Bigr)
= x + \frac13 - \frac{1}{9x} + o\Bigl(\frac1x\Bigr):
\]
assíntota $y = x + \frac13$, curva abaixo dela perto de $+\infty$. Quando $x
\to -\infty$,
o mesmo cálculo continua válido (a raiz cúbica está definida para
todos os reais, e $\frac1x \to 0$): mesma assíntota $y = x +
\frac13$, mas agora $-\frac{1}{9x} > 0$:
curva \emph{acima} da reta.
\end{solution}

\begin{solution}{exo:b1:taylor:8}
$u_n = \sqrt n\bigl(\sqrt{1 + \tfrac1n} - 1\bigr) = \sqrt
n\bigl(\frac{1}{2n} + o(\frac1n)\bigr) \sim \dfrac{1}{2\sqrt n}$.

$v_n = \ln\bigl(1 + \frac1n\bigr) \sim \dfrac 1n$.

$w_n$: com $h = \frac1n \to 0$, $\sin h - \tan h = \bigl(h -
\frac{h^3}{6}\bigr) - \bigl(h + \frac{h^3}{3}\bigr) + o(h^3) =
-\frac{h^3}{2} + o(h^3)$, logo $w_n \sim -\dfrac{1}{2n^3}$.
\end{solution}

\begin{solution}{exo:b1:taylor:9}
Taylor--Lagrange na ordem $1$ em torno de $x$, dos dois lados:
\[
f(x + h) = f(x) + h f'(x) + R_+,\quad
f(x - h) = f(x) - h f'(x) + R_-,
\qquad \abs{R_\pm} \leq \frac{h^2}{2} \sup \abs{f''} .
\]
Subtraindo: $f(x+h) - f(x-h) = 2h f'(x) + (R_+ - R_-)$, logo
\[
\abs{f'(x)} \leq \frac{\abs{f(x+h) - f(x-h)}}{2h} +
\frac{h}{2}\sup_{\intcc{x-h}{x+h}}\abs{f''} .
\]
Com as cotas globais: $\abs{f'(x)} \leq \frac{M_0}{h} + \frac{M_2
h}{2}$ para todo $h > 0$. O lado direito é
minimizado em $h =
\sqrt{2M_0/M_2}$ (\omterm{def:b1:derivative:def}{derivada} nula), com valor $\sqrt{2M_0M_2}$
--- logo $\abs{f'} \leq \sqrt{2M_0M_2}$. (Se $M_2 = 0$, faça $h \to
\infty$: $f' = 0$, o que é
coerente.)
\end{solution}

\begin{solution}{exo:b1:taylor:10}
Em $\intoo{0}{\pi}$: $0 < \sin u < u$, logo $(u_n)$ é estritamente
decrescente e positiva, portanto convergente; o limite é um ponto fixo
de $\sin$ em $\intcc{0}{\pi}$, e $\sin \ell = \ell$ força $\ell =
0$ (pois $\sin x < x$ para $x > 0$).
\begin{enumerate}
  \item Usando $\sin u = u - \frac{u^3}{6} + o(u^3)$ quando $u \to 0$:
        \[
        v_{n+1} - v_n = \frac{1}{\sin^2 u_n} - \frac{1}{u_n^2}
        = \frac{1}{u_n^2}\Bigl(\Bigl(1 - \frac{u_n^2}{6} +
        o(u_n^2)\Bigr)^{\!-2} - 1\Bigr)
        = \frac{1}{u_n^2}\Bigl(\frac{u_n^2}{3} + o(u_n^2)\Bigr)
        \longrightarrow \frac13 .
        \]
  \item Por \cref{exo:b1:seq:10}~(3) (Cesàro para diferenças),
        $\frac{v_n}{n} \to \frac13$, isto é, $v_n \sim \frac n3$, isto é,
        $u_n^2 \sim \frac 3n$: como $u_n > 0$,
        \[
        u_n \sim \sqrt{\frac{3}{n}} .
        \]
\end{enumerate}
\end{solution}

\begin{solution}{exo:b1:taylor:11}
Denominadores comuns. Primeiro limite:
\[
\frac{1}{x^2} - \frac{1}{\sin^2 x}
= \frac{\sin^2 x - x^2}{x^2\sin^2 x},
\qquad
\sin^2 x = \Bigl(x - \frac{x^3}{6} + o(x^4)\Bigr)^{\!2}
= x^2 - \frac{x^4}{3} + o(x^5) ,
\]
logo o numerador é $-\frac{x^4}{3} + o(x^4)$ enquanto o
denominador é $\sim x^4$: o limite é $-\dfrac13$.

Segundo: $\dfrac1x - \dfrac{1}{\ln(1+x)} = \dfrac{\ln(1+x) -
x}{x\ln(1+x)}$; o numerador é $-\frac{x^2}{2} + o(x^2)$, o
denominador $x\bigl(x + o(x)\bigr) \sim x^2$: o limite é
$-\dfrac12$.
\end{solution}

\begin{solution}{exo:b1:taylor:12}
Em $I_k = \intoo{k\pi - \frac\pi2}{k\pi + \frac\pi2}$, a
função $g(x) = \tan x - x$ tem \omterm{def:b1:derivative:def}{derivada} $\tan^2 x \geq 0$,
que se anula apenas no único ponto $k\pi$: $g$ é estritamente
crescente em $I_k$ (\cref{cor:b1:derivative:monotone}~(2)),
com limites $-\infty$ e $+\infty$ nas extremidades: exatamente um
zero $x_k$. Para $k \geq 1$, $g(k\pi) = -k\pi < 0$, logo $x_k \in
\intoo{k\pi}{k\pi + \frac\pi2}$: escreva $x_k = k\pi + \frac\pi2 -
\varepsilon_k$ com $\varepsilon_k \in \intoo{0}{\frac\pi2}$.
Então
\[
x_k = \tan x_k = \tan\Bigl(\frac\pi2 - \varepsilon_k\Bigr)
= \frac{1}{\tan\varepsilon_k}
\quad\Longrightarrow\quad
\varepsilon_k = \arctan\frac{1}{x_k} ,
\]
usando $\tan\varepsilon_k = \frac{1}{x_k}$ e $\varepsilon_k \in
\intoo{0}{\frac\pi2}$. Como $x_k \geq k\pi \to \infty$:
$\varepsilon_k \to 0$, e
\[
\varepsilon_k = \arctan\frac{1}{x_k} = \frac{1}{x_k} +
O\Bigl(\frac{1}{x_k^3}\Bigr)
= \frac{1}{k\pi + O(1)} + O\Bigl(\frac{1}{k^3}\Bigr)
= \frac{1}{k\pi} + O\Bigl(\frac{1}{k^2}\Bigr) ,
\]
donde $x_k = k\pi + \frac\pi2 - \frac{1}{k\pi} +
o\bigl(\frac1k\bigr)$.
\end{solution}

\begin{solution}{pb:b1:taylor:1}
\textbf{1.} $S_{2n+1} - S_{2n-1} = a_{2n} - a_{2n+1} \geq 0$ e
$S_{2n+2} - S_{2n} = a_{2n+2} - a_{2n+1} \leq 0$, enquanto $S_{2n}
- S_{2n+1} = a_{2n+1} \to 0$: as sequências $(S_{2n+1})$,
$(S_{2n})$ são adjacentes e convergem para um $S$ comum
(\cref{thm:b1:seq:adjacent}) com $S_{2n+1} \leq S \leq S_{2n}$.
Para $n$ par: $S_{n+1} \leq S \leq S_n$ dá $-a_{n+1} \leq S -
S_n \leq 0$; para $n$ ímpar: $0 \leq S - S_n \leq a_{n+1}$. Nos dois
casos $\abs{S - S_n} \leq a_{n+1}$ e $S - S_n$ tem o sinal de
$(-1)^{n+1}$, o primeiro termo omitido. O decrescimento estrito torna
estrita cada desigualdade exibida, em particular $0 < \abs{S -
S_n} < a_{n+1}$.

\textbf{2.} $a_k = \frac{1}{k!}$ decresce estritamente para $0$:
a questão 1 aplica-se. Taylor--Lagrange
(\cref{thm:b1:taylor:lagrange}) para $\exp$ entre $-1$ e $0$:
$\abs{\eu^{-1} - T_n} \leq \frac{1}{(n+1)!}$ (a \omterm{def:b1:derivative:def}{derivada}
$\eu^t$ é $\leq 1$ ali), logo $T_n \to \eu^{-1}$, e o limite
$S$ da questão 1 \emph{é} $\eu^{-1}$, com as cotas estritas
$0 < \abs{\eu^{-1} - T_n} < \frac{1}{(n+1)!}$.

\textbf{3.} Integrando a identidade em $\intcc{0}{1}$: o
lado esquerdo é $\arctan 1 = \frac\pi4$ (teorema fundamental), o
$k$-ésimo termo dá $\frac{(-1)^k}{2k+1}$, e
\[
\rho_n = (-1)^{n+1}\int_0^1 \frac{t^{2n+2}}{1+t^2}\dd t,
\qquad
\abs{\rho_n} \leq \int_0^1 t^{2n+2}\dd t = \frac{1}{2n+3} .
\]

\textbf{4.} O erro em $\pi$ é $4\abs{\rho_n} \leq
\frac{4}{2n+3}$: ficar abaixo de $10^{-6}$ exige
$2n + 3 > 4\cdot10^6$, cerca de dois milhões de termos. Já $4S_4 = 4\bigl(1 - \frac13 +
\frac15 - \frac17 + \frac19\bigr) = 4 \times 0.834921 =
3.339683$, a quase $0.2$ de
$\pi$: cinco termos, nem sequer um dígito.

\textbf{5.} Integrando em $\intcc{0}{1}$: $\ln 2 =
\sum_{k=1}^{n} \frac{(-1)^{k-1}}{k} + (-1)^n R_n$ com $R_n =
\int_0^1 \frac{t^n}{1+t}\dd t$; de $\frac12 \leq \frac{1}{1+t}
\leq 1$: $\frac{1}{2(n+1)} \leq R_n \leq \frac{1}{n+1}$. O
erro fica preso entre dois múltiplos de $\frac1n$: seis
decimais custam cerca de um milhão de termos.

\textbf{6.} Integrando de $0$ a $x$: $\frac12\ln\frac{1 +
x}{1 - x} = \sum_{k=0}^{n} \frac{x^{2k+1}}{2k+1} + \int_0^x
\frac{t^{2n+2}}{1-t^2}\dd t$. Em $x = \frac13$: $\frac{1 +
1/3}{1 - 1/3} = 2$, e em
$\intcc{0}{\frac13}$, $\frac{1}{1 -
t^2} \leq \frac98$:
\[
\ln 2 = 2\sum_{k=0}^{n} \frac{(1/3)^{2k+1}}{2k+1} +
\tilde\rho_n,
\qquad
0 < \tilde\rho_n \leq \frac{9}{4}\cdot
\frac{(1/3)^{2n+3}}{2n+3} :
\]
cada termo adicional \omterm{def:b1:arith:divides}{divide} o erro por cerca de $9$.

\textbf{7.} Indução: para $n = 1$: $1 - \frac12 = \frac12 = H_2
- H_1$. Passo:
\[
H_{2n+2} - H_{n+1} = (H_{2n} - H_n) + \frac{1}{2n+1} +
\frac{1}{2n+2} - \frac{1}{n+1}
= (H_{2n} - H_n) + \frac{1}{2n+1} - \frac{1}{2n+2},
\]
que é exatamente o incremento da soma alternada. E
$H_{2n} - H_n = \sum_{k=1}^{n}\frac{1}{n+k}$ é a soma de Riemann do \cref{ex:b1:integration:riemannexample}, que converge para $\ln
2$: as somas parciais pares da Via 1 \emph{são} as \omterm{thm:b1:integration:riemann}{somas de Riemann}
da Via 3.

\textbf{8.} $\sum_{k=1}^{6}\frac{(-1)^{k-1}}{k} = 0.61667$,
erro $0.0765$; a Via 2 em $n = 5$ dá $0.6931471$ com erro
$\leq \frac94\cdot\frac{(1/3)^{13}}{13} = 1.1\cdot10^{-7}$. A
razão: a Via 1 avalia a série do logaritmo no ponto de \omterm{def:b1:topology:closure}{fronteira}
$x = 1$, onde os termos decaem como $\frac1k$; a Via 2
avalia em $x = \frac13$, bem no \omterm{def:b1:topology:closure}{interior}, onde cada termo carrega
um fator novo $\frac19$.

\textbf{9.} Dez decimais: queremos $\tilde\rho_n \leq
5\cdot10^{-11}$. Em $n = 10$: $\frac94 \cdot
\frac{(1/3)^{23}}{23} = \frac94\cdot\frac{1.06\cdot10^{-11}}{23}
\approx 1.0\cdot10^{-12} < 5\cdot10^{-11}$: onze
termos bastam.

\textbf{10.} $(5+\iu)^2 = 24 + 10\iu$, e depois $(5+\iu)^4 = (24 +
10\iu)^2 = 476 + 480\iu$; e $2(1+\iu)(239+\iu) = 2(238 +
240\iu) = 476 + 480\iu$: iguais. Argumentos:
$\arg(5 + \iu) =
\arctan\frac15$, logo o lado esquerdo tem argumento
$4\arctan\frac15 \approx 0.79 \in \intoo{0}{\pi}$; o lado
direito tem argumento $\frac\pi4 + \arctan\frac{1}{239} \in
\intoo{0}{\pi}$. Dois números complexos iguais com
argumentos no mesmo \omterm{prop:b1:reals:intervals}{intervalo} de comprimento $< 2\pi$:
\[
4\arctan\frac15 = \frac\pi4 + \arctan\frac{1}{239} ,
\]
que é a fórmula de Machin.

\textbf{11.} Integre $\frac{1}{1+t^2} = \sum_{k=0}^n (-1)^k
t^{2k} + \frac{(-1)^{n+1}t^{2n+2}}{1+t^2}$ de $0$ a $x$:
\[
\arctan x = \sum_{k=0}^{n}\frac{(-1)^k x^{2k+1}}{2k+1} +
r_n(x),
\qquad
\abs{r_n(x)} \leq \int_0^x t^{2n+2}\dd t =
\frac{x^{2n+3}}{2n+3} .
\]

\textbf{12.} Erros: $16\,\frac{(1/5)^{11}}{11} = 3.0\cdot
10^{-8}$ e $4\,\frac{(1/239)^5}{5} \approx 1\cdot10^{-12}$:
total $< 5\cdot10^{-8}$. A soma exibida vale
$3.14159268\dots$, logo $\pi = 3.1415926\dots$ certificado com
$5\cdot10^{-8}$: sete decimais a partir de seis termos (cinco em
$\frac15$, dois em $\frac1{239}$ contando com generosidade).

\textbf{13.} Leibniz: erro $\sim \frac1n$, de modo que cada novo dígito
multiplica o trabalho por dez. Dalzell
(\cref{pb:b1:integration:1}, questão 22): erro $4^{1-5m}$,
cerca de três dígitos por passo, cada passo com um \omterm{def:b1:poly:def}{polinômio} mais
pesado. Machin: razão de erro $\frac{1}{25}$ por termo, cerca de $1.4$ dígitos
por termo, cada termo uma divisão. A moral: o resto de uma
expansão de tipo geométrico escala como $x^{2n}$, logo tornar $x$
pequeno compra dígitos a um custo \emph{fixo} por termo --- a
identidade complexa de Machin é precisamente uma máquina para
encolher $x$.

\textbf{14.} $a_k = \frac{1}{(2k)!}$ decresce estritamente para $0$;
pela questão 1 e pela \cref{prop:b1:taylor:standard} (cota de Lagrange
como na questão 2), $\sum_{k \leq n}\frac{(-1)^k}{(2k)!}
\to \cos 1$ com o enquadramento \emph{estrito} $0 < \bigl|\cos 1 -
S'_n\bigr| < \frac{1}{(2n+2)!}$.
Suponha $\cos 1 = \frac pq$ e tome $2n \geq q$: então $(2n)!\,S'_n = \sum_{k\leq n} (-1)^k
\frac{(2n)!}{(2k)!} \in \Z$ e $(2n)!\,\frac pq \in \Z$, ao passo que
\[
0 < \Bigl|(2n)!\cos 1 - (2n)!S'_n\Bigr| <
\frac{(2n)!}{(2n+2)!} = \frac{1}{(2n+1)(2n+2)} < 1 :
\]
um inteiro não nulo de valor absoluto $< 1$. Contradição: $\cos
1 \notin \Q$.

\textbf{15.} Identicamente com $a_k = \frac{1}{(2k+1)!}$,
multiplicando por $(2n+1)!$ com $2n + 1 \geq q$: $\sin 1 \notin
\Q$. E no entanto $\cos^2 1 + \sin^2 1 = 1 \in \Q$:
produtos e somas de irracionais podem ser racionais --- a
irracionalidade não atravessa de graça nenhuma operação algébrica.

\textbf{16.} Majoração da cauda: para $m > n$,
\[
\sum_{k=n+1}^{m} \frac{1}{(2k)!}
\leq \frac{1}{(2n+2)!}\Bigl(1 + \frac12 + \frac14 +
\dots\Bigr) \leq \frac{2}{(2n+2)!} ,
\]
pois cada razão sucessiva é $\frac{1}{(2k+1)(2k+2)} \leq
\frac12$; a cauda é positiva (o seu primeiro
termo é). Logo $0 <
\cosh 1 - \sum_{k\leq n}\frac{1}{(2k)!} < \frac{2}{(2n+2)!}$, e multiplicar por $(2n)!$ com $2n \geq q$ prende de
novo um inteiro não nulo em $\intoo{0}{1}$: $\cosh 1 \notin \Q$.

\textbf{17.} $\cos\frac1m = \sum_k
\frac{(-1)^k}{m^{2k}(2k)!}$: os termos decrescem estritamente para
zero, e $m^{2n}(2n)!\cdot\frac{1}{m^{2k}(2k)!} =
m^{2(n-k)}\frac{(2n)!}{(2k)!} \in \Z$ para $k \leq n$. Se
$\cos\frac1m = \frac pq$, multiplique o enquadramento estrito por
$q\,m^{2n}(2n)!$: o erro é majorado por
$\frac{q}{m^2(2n+1)(2n+2)} < 1$ para $n$ grande: contradição.
Para $\frac ab$ com $a \geq 2$: eliminar os denominadores multiplica
a cauda por $b^{2n}(2n)!$, mas o primeiro termo omitido é
$\frac{a^{2n+2}}{b^{2n+2}(2n+2)!}$, e o produto
$\frac{a^{2n+2}}{b^2(2n+1)(2n+2)}$ \emph{explode}: o
numerador $a^{2n+2}$ deixa de ser eliminado, e a armadilha emperra.
(O resultado continua verdadeiro --- por uma maquinaria ao estilo de
Niven, não por esta.)

\textbf{18.} Em $t = x$ todo termo de $g$ se anula exceto $f(x)
- f(x) = 0$:
$g(x) = 0$; $A$ é escolhido de modo que $g(a) = 0$.
Derivando, a soma telescopa:
\[
g'(t) = -\frac{f^{(n+1)}(t)}{n!}(x - t)^n +
A\,\frac{(x-t)^n}{n!}
= \frac{(x-t)^n}{n!}\bigl(A - f^{(n+1)}(t)\bigr) .
\]
Rolle no segmento de $a$ a $x$ dá $c$ estritamente entre eles
com $g'(c) = 0$; como $(x - c)^n \neq 0$: $A =
f^{(n+1)}(c)$. Desdobrando $g(a) = 0$ obtém-se a igualdade
de Taylor com resto $\frac{f^{(n+1)}(c)}{(n+1)!}(x-a)^{n+1}$.

\textbf{19.} Para $x > 0$, o resto é $\frac{\eu^{c}}
{(n+1)!}x^{n+1} > 0$: a exponencial supera cada
um dos seus \omterm{def:b1:poly:def}{polinômios} de Taylor, estritamente, em toda ordem. Para
$x < 0$ o sinal do resto é o de $x^{n+1}$: $\eu^x$ está
\emph{acima} do \omterm{def:b1:poly:def}{polinômio} para $n$ ímpar e \emph{abaixo} dele para
$n$ par --- lados alternados, como os gráficos de $1 + x$ e $1
+ x + \frac{x^2}{2}$
contra $\eu^x$ já mostram.

\textbf{20.} Erro verdadeiro: $\sin 0.5 - 0.4791667 =
2.59\cdot10^{-4}$, contra a cota $\frac{0.5^5}{120} =
2.60\cdot10^{-4}$: quase
atingida (o termo seguinte domina a cauda). Young: para limites e
análise local, em que nenhuma constante é necessária. Lagrange:
para decimais certificadas. Alternada: quando aplicável, a mesma
cota \emph{mais} a direção do erro --- a melhor das três, mas a
mais rara.

\textbf{21.} \omterm{def:b1:continuity:continuous}{Continuidade} em $0$: com $u = \frac{1}{x^2} \to
+\infty$, $f(x) = \eu^{-u} \to 0 = f(0)$. \omterm{def:b1:derivative:def}{Derivada} em $0$:
$\bigl|\frac{f(h)}{h}\bigr| = \sqrt u\,\eu^{-u} \to 0$
(\cref{prop:b1:functions:powerrules}): $f'(0) = 0$. Para $x \neq
0$, $f'(x) = \frac{2}{x^3}\eu^{-1/x^2}$: a forma
$P_1\bigl(\frac1x\bigr)\eu^{-1/x^2}$ com $P_1(X) = 2X^3$;
por indução, derivar $P_k(\frac1x)\eu^{-1/x^2}$ dá
$P_{k+1}(X) = 2X^3 P_k(X) - X^2 P_k'(X)$, um \omterm{def:b1:poly:def}{polinômio}. Então
\[
\frac{f^{(k)}(h) - 0}{h} = v\,P_k(v)\,\eu^{-v^2}
\Big|_{v = 1/h} \longrightarrow 0
\]
(\omterm{def:b1:poly:def}{polinômio} contra $\eu^{-v^2}$, comparação de crescimento em
$\pm\infty$): por indução $f^{(k)}(0) = 0$ para todo $k$. Todos os
\omterm{def:b1:poly:def}{polinômios} de Taylor de $f$ em $0$ se anulam, e no entanto $f > 0$
fora de $0$: Taylor--Young é exato em toda ordem e cego para além do
germe. Uma expansão é informação apenas local.

\textbf{22.} Pela questão 2, $0 < \abs{\eu^{-1} - T_n} <
\frac{1}{(n+1)!}$, estritamente. Se $\eu^{-1} = \frac pq$, tome $n
\geq q$
e multiplique por $n!$: $n!\,T_n \in \Z$ e $n!\frac pq
\in \Z$, de modo que um inteiro não
nulo tem valor absoluto $<
\frac{n!}{(n+1)!} = \frac{1}{n+1} < 1$: contradição. Logo
$\eu^{-1} \notin \Q$, e $\eu = \frac{1}{\eu^{-1}}$ também é
irracional. As três demonstrações: \omterm{thm:b1:seq:adjacent}{sequências adjacentes} espremendo
$q!\,\eu$ (\cref{exo:b1:seq:9}); a recorrência \omterm{thm:b1:integration:def}{integral} $A_n =
\eu - nA_{n-1}$ (\cref{pb:b1:integration:1}); o
enquadramento alternado (aqui). Uma armadilha, três certificados de
pequenez.

\textbf{23.} Agrupe as somas parciais aos pares: $\sum_{k=1}^{2n}
(-1)^k b_k = E_n - O_n$ com $E_n = \sum_{j=1}^{n}
\frac{1}{4j^2}$,
limitada (pela cota telescópica do \cref{ex:b1:seq:basel}, $E_n \leq \frac12$), e $O_n =
\sum_{j=1}^{n}\frac{1}{2j-1} \geq \frac12 H_n \to +\infty$
(\cref{exo:b1:seq:5}): as somas parciais tendem a $-\infty$.
A monotonia foi usada na questão 1 exatamente onde $S_{2n+1} -
S_{2n-1} = a_{2n} - a_{2n+1}$ precisava de
um sinal: sem decrescimento, as \omterm{def:b1:seq:subsequence}{subsequências} par e ímpar não
precisam ser monótonas, e a adjacência desaba.

\textbf{24.} $(2+\iu)(3+\iu) = 5 + 5\iu = 5(1+\iu)$; tomando
argumentos (todos em $\intoo{0}{\frac\pi2}$): $\arctan\frac12 +
\arctan\frac13 = \frac\pi4$. Custo em série para seis
decimais: erro $\leq 4\bigl(\frac{(1/2)^{2n+3}}{2n+3} +
\frac{(1/3)^{2n+3}}{2n+3}\bigr)$, que em $n = 10$ vale
$\approx 2\cdot10^{-8} < 5\cdot10^{-7}$: onze termos. Classificação:
melhor do que Leibniz por uma margem exponencial, atrás de Machin
(cujo ponto dominante $\frac15$ é menor do que $\frac12$): cerca de
$0.6$ dígitos por termo contra os $1.4$ de Machin.

\textbf{25.} (i) Para $a_k \to 0$ decrescente, as somas parciais alternadas
convergem com $\abs{S - S_n} \leq a_{n+1}$ e o erro carrega o sinal do primeiro termo
omitido. (ii) Uma identidade finita com resto explícito pode ser
\emph{avaliada e majorada num ponto escolhido}, enquanto um
enunciado de limite só promete proximidade eventual --- a
certificação exige a primeira. (iii) Extraído: $\pi$ com $5\cdot10^{-8}$ por
Machin, $\ln 2$ com dez decimais pela série em $\frac13$, a
irracionalidade de $\cos 1$, $\sin 1$, $\cosh 1$, $\cos\frac1m$ e $\eu^{-1}$, a
igualdade de Taylor--Lagrange e o aviso da função plana. (iv)
O \cref{ch:b1:series} transforma a questão 1 no teste das séries
alternadas, separa a convergência absoluta da condicional,
e o seu problema de fim de semana encena o drama do rearranjo,
para o qual a série harmônica alternada da Via 1 é a
testemunha principal.
\end{solution}
